(1) 在 $x>0$ 内定态薛定谔方程为 $-\hbar^2\psi''/(2m)+mgx\psi=E\psi$。无穷硬壁给 $\psi(0)=0$，同时要求 $\psi(x)$ 在 $x\to\infty$ 衰减并平方可积，因此只允许离散束缚能级。
(2) 取 $\ell^3=\hbar^2/(2m^2g)$，定义 $z=x/\ell-E/(mg\ell)$，代入后得到 $\psi_{zz}-z\psi=0$，即Airy微分方程。一般解是 $c_1\operatorname{Ai}(z)+c_2\operatorname{Bi}(z)$；因 $\operatorname{Bi}(z)$ 在正无穷增长，必须令 $c_2=0$，归一化由 $\int_0^\infty|\psi|^2dx=1$ 决定。
(3) 若 $\operatorname{Ai}(-a_n)=0$，其中 $a_1\simeq2.338107$、$a_2\simeq4.087949$ 为前两个正数，硬壁边界给 $E_n=mg\ell a_n$，且 $\psi_n(x)=C_n\operatorname{Ai}(x/\ell-a_n)$。第一激发间距为 $\Delta E_{21}=mg\ell(a_2-a_1)>0$，不等间距体现线性势的非简谐谱。
(4) 对线性势 $V=mgx$，量子维里定理给 $2\langle K\rangle=\langle xV'(x)\rangle=mg\langle x\rangle$，同时 $E=\langle K\rangle+mg\langle x\rangle=3mg\langle x\rangle/2$。因此 $\langle x\rangle_n=2E_n/(3mg)=2a_n\ell/3$。因为 $\ell\propto g^{-1/3}$，能量 $E_n\propto g^{2/3}$。
